\documentclass[10pt,letterpaper]{article}

\usepackage[margin=1in]{geometry}
\usepackage{graphicx}
\usepackage{xcolor}
\usepackage{amsmath}
\usepackage{amssymb}
\usepackage{booktabs}
\usepackage{caption}
\usepackage{subcaption}
\usepackage[hyphens]{url}
\usepackage[hidelinks]{hyperref}
\usepackage{authblk}
\usepackage{lineno}
\usepackage{titlesec}

\titleformat{\section}{\bfseries\normalsize}{\thesection}{0.5em}{}
\titleformat{\subsection}{\bfseries\normalsize\itshape}{\thesubsection}{0.5em}{}

% Round 9 P2: "Fig. N." bold caption prefix with period (matches reference bioRxiv style)
\renewcommand{\figurename}{Fig.}
\captionsetup{labelfont=bf,labelsep=period}

% Round 9 P5: ABSTRACT uppercase bold heading (matches reference bioRxiv style)
\renewcommand{\abstractname}{\textbf{ABSTRACT}}

\title{Aether3D: A Reproducible Benchmark Contract for a Planned Continuous 3D Spatial-Omics Reconstruction Claim}

\author[1]{Anonymous Authors}
\affil[1]{Author affiliations to be added before submission}

\date{\today}

\newcommand{\statusbanner}{%
  \begin{center}\colorbox{yellow!30}{\parbox{0.9\textwidth}{%
  \textbf{DRAFT --- SYNTHETIC PLACEHOLDER DATA.} Every numerical result and
  figure in this manuscript is generated from the synthetic benchmark JSONs
  under \path{aether-3d/results/benchmark/}. No claim has graduated to
  \emph{paper-ready} status in the project claim ledger. Real-data
  promotion is gated by the acceptance matrix.
  }}\end{center}}

\begin{document}
\maketitle
\linenumbers
\statusbanner


\begin{abstract}
Serial-section tissue imaging produces discrete 2D slices whose
inter-section geometry is often hypothesized to benefit from continuous
3D reconstruction. This draft now treats that positive reconstruction
claim as \emph{planned (contradicted)} rather than established: the
project claim ledger records that current real leave-one-out benchmarks
on two serial datasets favor simple 2.5D / linear-interpolation baselines
over the continuous adapter on the primary Moran's-I, Betti-0, and
sliced-Wasserstein metrics. We therefore present \textbf{\textsc{Aether3D}}
as a reproducible benchmark contract and repair target, not as a validated
advantage result. The contract wraps reconstruction adapters under a
single virtual-slice holdout protocol and reports Chamfer distance,
coordinate RMSE, sliced 2D Wasserstein, Moran's I agreement, domain ARI,
domain NMI, cell-type proportion Spearman, and Betti-0 topology stability.
Flow-divergence and velocity-anisotropy diagnostics are computed only
when a velocity field is available. All positive biological,
cell-fidelity, and continuous-over-2.5D performance claims remain gated by
the claim ledger and cannot graduate until repaired benchmarks reverse the
current contradictory evidence.
\end{abstract}

\section{Introduction}
\label{sec:intro}

Three-dimensional tissue context cannot be recovered from 2D imaging
without an explicit reconstruction step that addresses two coupled
problems: (i)~aligning discrete sections in a shared 3D coordinate
system; (ii)~estimating the cellular and molecular state at intermediate
depths the imager did not capture. The local reference stack now includes
multi-slice integration, alignment, and platform-benchmark papers that make
section spacing, missing sections, and modality-specific resolution explicit
sources of variation~\cite{Dong2025MultisliceBenchmark,
ClusteringIntegrationBenchmark2024,You2024SequencingST,Cervilla2026SixCancerPlatforms}.
Prior art spans alignment-only
methods~\cite{Zeira2022PASTE,Liu2023PASTE2,Zhou2023STAligner,Jones2023GPSA,Clifton2023STalign},
2.5D virtual-slice interpolators, and continuous-field
generative and multi-slice 3D-reconstruction models~\cite{Wang2023STitch3D};
we cite these as motivation rather than re-using their
methodology, with the baseline reference treated as an aspirational large-scale prior-art
comparison rather than a reproduced result~\cite{Yang2026DeepSpatial}.
\textsc{Aether3D} builds on entropic optimal
transport~\cite{Cuturi2013Sinkhorn} for slice coupling, flow
matching~\cite{Lipman2023FlowMatching} for the velocity field, and
adaptive normalization~\cite{Peebles2023DiT} for conditioning on
heterogeneous cell-state tokens. Downstream analyses use spatial
autocorrelation~\cite{Moran1950Autocorrelation}, adjusted Rand
similarity~\cite{Hubert1985Adjusted}, and persistent-homology
descriptors~\cite{Edelsbrunner2010PersistentHomology}; spatial-domain
application pressure is represented by CellCharter-style niche discovery,
which remains cited rather than integrated in this scaffold~\cite{Varrone2024CellCharter}.
The data plumbing
follows standard single-cell tooling~\cite{Wolf2018Scanpy}.

\textsc{Aether3D} contributes three components: (i)~a uniform
volume-adapter contract that lets methods be compared on identical
inputs under the held-out audit boundary;
(ii)~an eight-metric per-virtual-slice suite (Section~\ref{sec:metrics})
that reports geometry, molecular, and topology fidelity in one pass;
(iii)~an explicit $\lambda_{\text{class}}$ knob with a reproducible ablation
grid that examines the shipped $\alpha_{\text{spatial}} = 0.5$,
$\lambda_{\text{class}} = 10$ defaults.

Per the project audit boundary, no prose, class name, or figure layout
from any prior 3D reconstruction work is reused; the leakage scanner
enforces this.

\section{Method}
\label{sec:method}

\subsection{Problem formulation}
\label{sec:problem-formulation}

A serial-section spatial-omics volume is a tuple
$\mathcal{V} = \{(\mathbf{X}_k, \mathbf{S}_k, z_k)\}_{k=1}^K$ where for
each section $k$, $\mathbf{X}_k \in \mathbb{R}^{N_k \times G}$ is the
per-cell expression, $\mathbf{S}_k \in \mathbb{R}^{N_k \times 2}$ are
the 2D centroid coordinates, and $z_k \in \mathbb{R}$ is the physical
depth in micrometers. A 3D reconstruction adapter
$g_\phi\!:\!\mathcal{V} \mapsto \tilde{\mathcal{V}}(z)$ returns a
continuous family of synthetic slices at any query depth $z \in
[\min_k z_k, \max_k z_k]$. The audit boundary requires that
$g_\phi$ never observes a held-out slice $k^\star$ during the call
that asks for $\tilde{\mathcal{V}}(z_{k^\star})$.

\subsection{Architecture: volume-adapter contract + hybrid UOT coupling}
\label{sec:contract}

Each method we compare implements a uniform \texttt{VolumeBaseAdapter}
subclassing pattern. The visible input is a list of AnnData slices with
\texttt{.obsm["spatial"]} 2D coords and a per-slice physical-$z$ value;
held-out slices are removed from the adapter-visible input before any
call. Per-result provenance records device, runtime, peak memory,
torch version, CUDA version, hostname, and git SHA. Six test files
under \texttt{aether-3d/tests/benchmarks/} independently verify the
audit boundary across the sweep harnesses.

The 2.5D-stacking baseline (\texttt{linear-interp}) is a plain
flanking-slice linear interpolation; the continuous-3D path uses the
hybrid UOT coupling described in Section~\ref{sec:uot} together with a
flow-matching velocity field and an ODE-integrated density-preserving
sampler. The architecture is described once via the adapter contract so
new methods plug in without changing the runner.

\subsection{Hybrid UOT coupling}
\label{sec:uot}

The cell-to-cell coupling cost between adjacent slices combines spatial
distance, gene cosine distance, and a same-cell-type penalty. Let
$\hat{D}^{\text{spat}}$ and $\hat{D}^{\text{gene}}$ be the spatial
(Euclidean) and gene (cosine) distance matrices, each max-normalized to
$[0, 1]$, and let $P^{\text{cell}}_{ij} = \mathrm{clip}(1 - \langle
\mathbf{c}_i, \mathbf{c}_j \rangle, 0, 1)$ be the one-hot cell-type
mismatch between cells $i$ and $j$:
\[
  C_{ij} = \alpha_{\text{spatial}} \, \hat{D}^{\text{spat}}_{ij}
         + (1 - \alpha_{\text{spatial}}) \, \hat{D}^{\text{gene}}_{ij}
         + \lambda_{\text{class}} \, P^{\text{cell}}_{ij}.
\]
Both $\alpha_{\text{spatial}} \in [0, 1]$ and $\lambda_{\text{class}} \in
[0, \infty)$ are exposed as configuration parameters; the shipped
defaults are $\alpha_{\text{spatial}} = 0.5$ and $\lambda_{\text{class}} =
10$. The $(\alpha, \lambda)$ ablation grid in
Section~\ref{sec:res-ablation} quantifies the contribution of each
component and justifies those defaults.

\subsection{Training and inference protocol}
\label{sec:training-protocol}

The continuous-3D adapters expose a deterministic \texttt{fit / sample}
contract: \texttt{fit} consumes the visible slices plus their physical-$z$
labels and returns an adapter object; \texttt{sample(z)} returns the
reconstructed AnnData at any in-range depth. The
scaling harness (\texttt{aether\_3d.benchmarks.scaling}) records device,
torch version, CUDA version, hostname, peak GPU memory, runtime, and
git SHA per measurement, so a reviewer can re-run any single scaling
point. Downstream cellular-neighborhood analysis
(\texttt{aether\_3d.benchmarks.\allowbreak neighborhood}) computes radius-$R$
distance-band enrichment and z-axis depth-binned density profiles
around any query cell type, supporting the in-silico sectioning and
co-localization analyses described in Section~\ref{sec:results}.

\subsection{Metric suite}
\label{sec:metrics}

For each held-out virtual depth $z_t$ the reconstructed cells are
compared against the truth slice at that depth on eight metrics:
Chamfer distance and coordinate RMSE (geometry), sliced 2D Wasserstein
(distributional geometry), Moran's I agreement on top-$K$ HVGs
(molecular continuity), domain ARI and domain NMI via KMeans on the
latent (cluster identity), cell-type proportion Spearman (identity
composition), and Betti-0 stability via kNN union-find (topology).

\textbf{Metric definitions.}
\emph{Chamfer distance}: the mean of the per-direction nearest-neighbor
Euclidean distances between two point clouds $A$ and $B$,
$d_\text{Ch}(A,B)=\tfrac{1}{2}(\bar{d}_{A\to B}+\bar{d}_{B\to A})$
($\downarrow$ better).
\emph{Coordinate RMSE}: root-mean-squared nearest-neighbor distance
from the reconstruction to truth ($\downarrow$ better).
\emph{Sliced 2D Wasserstein}: approximation of the 2D Wasserstein
distance via mean 1D Wasserstein over $n=50$ random projections;
measures distributional shift in cell position ($\downarrow$ better).
\emph{Moran's I agreement}: Spearman correlation of per-gene
Moran's~$I$ statistics~\cite{Moran1950Autocorrelation} computed
independently on truth and reconstruction, evaluated over the top-$K$
(default $K=100$) highest-variance genes; measures whether the spatial
autocorrelation structure of gene expression is preserved
($\uparrow$ better).
\emph{Domain ARI / NMI}: KMeans labels fit to the pooled
truth-plus-reconstruction gene-expression matrix; ARI~\cite{Hubert1985Adjusted}
and NMI between the two halves measure cluster-identity agreement
($\uparrow$ better).
\emph{Cell-type proportion Spearman}: Spearman correlation of per-label
cell counts between truth and reconstruction ($\uparrow$ better).
\emph{Betti-0 stability}: fraction of $k$-NN connected components
preserved after swapping truth cells for reconstructed cells, computed
via union-find over the joint kNN
graph~\cite{Edelsbrunner2010PersistentHomology} ($\uparrow$ better).
Flow-divergence and velocity-anisotropy vector-field diagnostics are
computed by the same benchmark module when a velocity field is available
from a trained continuous-3D adapter; they are not applicable to the
NumPy-only 2.5D baselines.

\subsection{Leave-one-section-out evaluation protocol}
\label{sec:loso}

The primary evaluation uses a \emph{leave-one-section-out} (LOSO)
holdout: given $K$ aligned serial sections $\{(\mathbf{X}_k,
\mathbf{S}_k, z_k)\}_{k=1}^K$, section $k^\star$ is designated the
held-out truth; the adapter receives the remaining $K-1$ sections as
visible input and is asked to produce a virtual slice at depth
$z_{k^\star}$. The protocol iterates over all interior holdout
positions $k^\star \in \{2, \ldots, K-1\}$ and reports mean~$\pm$~std
of each metric across holdout rounds.

The audit boundary is enforced mechanically by the
\texttt{VolumeAdapterInput} contract
(Section~\ref{sec:contract}): \texttt{visible\_slices()} removes the
held-out section from the input list before the adapter's \texttt{fit}
call, and \texttt{truth\_slices()} exposes the truth only to the scorer.
An adapter that attempts to read beyond \texttt{visible\_slices()} has
no path to the truth expression matrix --- this is a closed-form
data-leakage guard rather than an honor system.
Per-holdout-slice scores are stored in a structured JSON keyed by
$(k^\star, \text{method})$ so that any single evaluation point can be
re-run independently with its stored provenance record (device, git SHA,
torch and CUDA versions, hostname, seed).
The aggregated means and standard deviations across all holdout rounds
constitute the primary comparison surface
(Table~\ref{tab:loso-gene}).

\subsection{Baselines}
\label{sec:baselines}

Three parameter-free 2.5D baselines and the \textsc{Aether3D}
continuous-3D model are wired through the same
\texttt{VolumeBaseAdapter} contract and scored under the identical LOSO
protocol.

\textbf{Nearest-section} (\texttt{nearest-slice}): the held-out depth
$z_{k^\star}$ is reconstructed by copying all cells from the single
closest visible section and overwriting their $z$ coordinate to the
target depth. This is the lowest possible 2.5D floor: it uses real
expression and real 2D positions with no interpolation.

\textbf{Linear interpolation} (\texttt{linear-interp}): each
virtual-depth cell is produced by linearly interpolating expression
vectors and 2D coordinates between the two bracketing visible sections.
Cells are paired by nearest-neighbour in 2D (the only correspondence
signal available without ground-truth cell tracks) and the interpolation
weight is $w = (z_{k^\star} - z_a)/(z_b - z_a)$. The pairing is
deterministic in cell content (ordered by \texttt{lexsort} on 2D
coordinates) so reported numbers are invariant to input row order.

\textbf{2.5D stacking} (\texttt{stacking-2.5d}): the canonical
clean-room reimplementation of the 2.5D virtual-slice stacking strategy.
A virtual section at depth $z$ is assembled by stacking
$k_a = \lfloor(1-w)\,n_a\rceil$ cells from the lower flanking section
and $k_b = \lfloor w\,n_b\rceil$ from the upper one, selected by an
evenly-spaced decimation over a lexsort-deterministic spatial ordering
so the subset is spatially representative and independent of row order.
Every selected cell carries its real expression vector and real 2D
position; its $z$ coordinate is overwritten to the target depth. There
is no learned model and no gene-wise smoothing, so any quality gap
between this baseline and a continuous method is attributable to the
method, not to artefacts on the baseline side.

\textbf{Aether continuous-3D} (\texttt{aether}): the
\textsc{Aether3D} reconstructor is wrapped in an \texttt{AetherAdapter}
that implements the identical \texttt{VolumeBaseAdapter} interface and
is scored under the same LOSO protocol as the three 2.5D baselines.
The adapter calls \texttt{AetherReconstructor.fit} on the visible
slices, then \texttt{reconstruct\_continuous\_volume} at the held-out
depth. Physical-$z$ remapping from the reconstructor's internal depth
index back onto the observed $\mu$m values is handled inside the
adapter, so the scorer always operates on the same physical-$\mu$m
coordinate system as the baselines. Results are placeholders pending a
validated run (Table~\ref{tab:loso-gene}); no performance claim over
baselines is asserted here.

\subsection{Fair z-axis: physical micrometres vs.\ synthetic uniform
  spacing}
\label{sec:fair-z}

Serial-section datasets carry \emph{physically irregular}
inter-section gaps: adjacent sections may be separated by anywhere from
tens to several hundred micrometres depending on cutting protocol,
tissue heterogeneity, and section loss. All four adapters in this
benchmark consume physical $z_k$ values in micrometres as stored in
\texttt{obs["z\_coord"]} rather than a synthetic integer index.

This distinction matters asymmetrically across method families. The
2.5D baselines (nearest-section, linear interpolation, stacking) are
\emph{scale-invariant}: their interpolation weight
$w = (z_{k^\star} - z_a)/(z_b - z_a)$ is a ratio that normalises out
any common rescaling of the $z$ axis.  A 2.5D method asked to
interpolate between $z_a = 0\,\mu$m and $z_b = 200\,\mu$m at
$z_{k^\star} = 100\,\mu$m yields $w = 0.5$, the same as if the
coordinates were integers $0$, $1$, $2$.

The \textsc{Aether3D} continuous-3D model is \emph{z-conditioned}: the
velocity field takes physical $z$ as an input to the adaptive
normalisation layer (Section~\ref{sec:uot}), and the ODE integration
time is derived from the physical spacing between slices. Evaluating
a $z$-conditioned model on a synthetic uniform-spacing approximation
introduces a systematic bias when real sections are irregularly spaced:
the model is trained and queried at wrong time intervals, inflating
reconstruction error independently of model quality.

Feeding physical $\mu$m coordinates to both 2.5D baselines (which are
unaffected by coordinate scale) and to \textsc{Aether3D} (which depends
on them) is therefore a methodological contribution to fair evaluation:
it ensures that any observed difference between the continuous-3D method
and the 2.5D floor reflects the model's ability to exploit inter-section
geometry rather than an artefact of coordinate normalisation. The
contract enforces this at ingestion time: the serial data loader raises
\texttt{ValueError} if \texttt{obs["z\_coord"]} is absent, preventing
silent fallback to synthetic uniform spacing.

\subsection{Falsifiability controls}
\label{sec:falsif}

The LOSO evaluation is accompanied by two synthetic controls that
verify whether the metric-plus-protocol combination can
\emph{detect a difference} between method families, independent of
any real-data claim. Both are generated from one parametric field in
which the \emph{same} cells are observed at every section, each tracing
a deterministic trajectory $\mathbf{p}(z) = \mathbf{p}_0 +
\mathbf{v}\,z + c\,\mathbf{d}\,z^2$ across $z$ (with $z$ centred so the
held-out interior section sits at $z=0$). A single curvature knob $c$
switches between the two regimes.

\textbf{Negative control --- regime where 2.5D is provably optimal
($c = 0$).}
With zero curvature the trajectories are linear, so the true held-out
midpoint section is \emph{exactly} the linear blend of its two flanking
sections (the $\mathbf{v}\,z$ term cancels at the centred midpoint).
Linear interpolation is therefore near-exact and a continuous model has
nothing to gain: a tie or loss for the learned adapter is the correct
outcome. A spurious win for the learned model in this regime would
indicate a calibration failure of the metrics rather than a genuine
advantage.

\textbf{Positive control --- regime where linear interpolation is
provably biased ($c > 0$).}
With nonzero curvature the quadratic bend makes the true midpoint differ
from the linear blend of the flanking sections by a fixed offset
$c\,\mathbf{d}$, so linear interpolation systematically misplaces cells
and is provably biased on Chamfer distance and coordinate RMSE.  A
continuous-flow model conditioned on physical $z$ has the
representational capacity to recover the curve; a scale-invariant 2.5D
interpolator does not.  This control establishes that the benchmark
\emph{can} detect a continuous-method advantage when one exists, without
asserting that \textsc{Aether3D} achieves it on real data (that claim
remains \texttt{planned} in the claim ledger; see
Table~\ref{tab:loso-gene}).

Both controls are generated deterministically and scored through the
identical adapter contract.  The \emph{validity} of the controls
themselves --- linear interpolation near-exact in the negative regime
and provably biased in the positive regime --- is asserted in the test
suite, so a regression in the metric or protocol implementation surfaces
immediately.

% ---------------------------------------------------------------------------
% Placeholder tables and figures (numbers TBD)
% ---------------------------------------------------------------------------

\subsection{Placeholder result tables and figures}
\label{sec:placeholders}

The following tables and figures are structural stubs. Captions describe
what each item will contain; all numerical entries are marked
\textbf{TBD} and will be filled once the claim ledger entries for the
relevant rows reach \texttt{validated} status with linked, reproducible
artifacts. The table and figure environments are kept here so the
reference labels are stable and the document compiles cleanly.

\begin{table}[ht]
\centering
\caption{LOSO gene-recovery table --- mean~$\pm$~std of each metric
across leave-one-section-out rounds, all adapters scored through the
same audited contract. Adapters: \texttt{nearest-slice},
\texttt{linear-interp}, \texttt{stacking-2.5d}, \texttt{aether}.
Metrics: Chamfer ($\downarrow$), coord.\ RMSE ($\downarrow$),
sliced Wasserstein ($\downarrow$), Moran's I agreement ($\uparrow$),
domain ARI ($\uparrow$), domain NMI ($\uparrow$), cell-type proportion
Spearman ($\uparrow$), Betti-0 ($\uparrow$).
\textbf{All entries are TBD --- pending validated run on real serial
data.}}
\label{tab:loso-gene}
\small
\begin{tabular}{lcccccccc}
\toprule
Method & Chamfer & RMSE & SW$_{2}$ & Moran's I & ARI & NMI & Spearman & Betti-0 \\
\midrule
nearest-slice  & TBD & TBD & TBD & TBD & TBD & TBD & TBD & TBD \\
linear-interp  & TBD & TBD & TBD & TBD & TBD & TBD & TBD & TBD \\
stacking-2.5d  & TBD & TBD & TBD & TBD & TBD & TBD & TBD & TBD \\
aether         & TBD & TBD & TBD & TBD & TBD & TBD & TBD & TBD \\
\bottomrule
\end{tabular}
\end{table}

\begin{table}[ht]
\centering
\caption{Per-region cell-type proportion Spearman --- Spearman
correlation of per-cell-type counts between reconstruction and truth,
reported separately for each annotated spatial region
(e.g., cortical layer, hypothalamic nucleus) on the holdout dataset.
Columns correspond to regions; rows to adapters.
\textbf{All entries are TBD --- pending validated run.}}
\label{tab:per-region-spearman}
\small
\begin{tabular}{lccc}
\toprule
Method & Region 1 & Region 2 & $\cdots$ \\
\midrule
nearest-slice  & TBD & TBD & $\cdots$ \\
linear-interp  & TBD & TBD & $\cdots$ \\
stacking-2.5d  & TBD & TBD & $\cdots$ \\
aether         & TBD & TBD & $\cdots$ \\
\bottomrule
\end{tabular}
\end{table}

\begin{table}[ht]
\centering
\caption{Cross-platform summary --- mean Chamfer distance and Moran's I
agreement on each validated dataset (e.g., MERFISH mouse hypothalamus,
BRCA IMC, Visium serial ladder), broken down by adapter. This table will
exist only when at least two distinct real-data platforms have completed
the LOSO protocol.
\textbf{All entries are TBD --- pending multi-platform validated runs.}}
\label{tab:cross-platform}
\small
\begin{tabular}{lccc}
\toprule
 & \multicolumn{2}{c}{Platform A} & Platform B \\
\cmidrule(lr){2-3}
Method & Chamfer & Moran's I & Chamfer \\
\midrule
nearest-slice  & TBD & TBD & TBD \\
linear-interp  & TBD & TBD & TBD \\
stacking-2.5d  & TBD & TBD & TBD \\
aether         & TBD & TBD & TBD \\
\bottomrule
\end{tabular}
\end{table}

\begin{figure}[ht]
  \centering
  \fbox{\parbox{0.9\textwidth}{\centering\vspace{2em}
  \textbf{[Figure stub: intermediate-depth montage]}\\[0.5em]
  Side-by-side comparison of truth vs.\ reconstructed virtual sections
  at three intermediate depths per holdout round. Columns: truth,
  nearest-slice, linear-interp, stacking-2.5d, aether. Rows: holdout
  depths $z_{k^\star}$ (e.g., $z_2$, $z_3$, $z_4$ for a 6-section
  stack). Cell positions shown as scatter; colour = cell type.\\[0.5em]
  \textbf{TBD --- pending validated run on real serial data.}
  \vspace{2em}}}
  \caption{Intermediate-depth reconstruction montage. Each panel shows
  cells coloured by annotated cell type at a held-out intermediate
  depth; the truth section (left) is compared to each adapter's
  reconstruction. \textbf{TBD --- pending validated run.}}
  \label{fig:depth-montage}
\end{figure}

\begin{figure}[ht]
  \centering
  \fbox{\parbox{0.9\textwidth}{\centering\vspace{2em}
  \textbf{[Figure stub: velocity/warp field]}\\[0.5em]
  Visualisation of the \textsc{Aether3D} learned velocity field over
  the 3D tissue volume. Arrows or streamlines represent the
  ODE-integrated flow from section $z_k$ to $z_{k+1}$; colour encodes
  velocity magnitude. Divergence and anisotropy diagnostics
  are displayed as scalar heat maps.\\[0.5em]
  \textbf{TBD --- pending trained model run with velocity-field output
  enabled.}
  \vspace{2em}}}
  \caption{Velocity/warp field from the continuous-3D adapter.
  Streamlines show the inferred cell-movement field between adjacent
  sections; scalar panels show flow divergence and velocity anisotropy.
  \textbf{TBD --- pending trained model run.}}
  \label{fig:velocity-field}
\end{figure}

\begin{figure}[ht]
  \centering
  \fbox{\parbox{0.9\textwidth}{\centering\vspace{2em}
  \textbf{[Figure stub: marker concordance]}\\[0.5em]
  Per-gene Pearson correlation (reconstruction vs.\ truth, spatially
  matched) for the top marker genes of each annotated cell type,
  shown as a dot plot or heatmap. Rows = genes; columns = adapters.
  Reference line at $r = 0$ marks the chance level.\\[0.5em]
  \textbf{TBD --- pending validated run.}
  \vspace{2em}}}
  \caption{Marker-gene expression concordance. Spatially-matched
  per-gene Pearson correlation between reconstruction and truth for
  top cell-type marker genes. \textbf{TBD --- pending validated run.}}
  \label{fig:marker-concordance}
\end{figure}

\section{Results}
\label{sec:results}

\subsection{Overview of \textsc{Aether3D}}
\label{sec:res-overview}

\textsc{Aether3D} has three independently-testable components.
\emph{(i)}~A uniform volume-adapter contract (Section~\ref{sec:contract})
runs the 3D-reconstruction methods on identical serial-section input with
held-out virtual slices removed before any adapter sees the data.
\emph{(ii)}~An eight-metric per-virtual-slice suite
(Section~\ref{sec:metrics}) reports geometry, distributional geometry,
molecular continuity, cluster identity, and topology fidelity in a single
pass. \emph{(iii)}~A scaling harness records device,
torch version, CUDA version, peak memory, and runtime per measurement.
The synthetic benchmarks below exercise each of these; real-data
graduation remains gated by the project claim ledger.

\subsection{Virtual-slice holdout (synthetic placeholder)}
\label{sec:res-holdout}

Figure~\ref{fig:holdout} reports the per-virtual-slice metrics for the
two always-available baselines (nearest-slice, linear-interp) on a
synthetic 5-slice stack with the middle slice held out. Nearest-slice
attains lower Chamfer distance and lower sliced 2D Wasserstein than
linear interpolation at this noise level, which has a simple
explanation: with $\sigma_{\text{spatial}} = 1.0$ the nearest
neighbor is closer to the truth than the linear blend of two flanking
slices. \emph{Numbers in this figure are synthetic placeholders; real-data
figures will replace them once a serial-section dataset graduates.}

\begin{figure}[t]
  \centering
  \includegraphics[width=\textwidth]{figures/fig_holdout_quartet.pdf}
  \caption{Virtual-slice holdout metrics for the two
  always-available baselines on a synthetic 5-slice stack with the
  middle slice held out. Panels \textbf{a}--\textbf{g} show seven of the
  eight per-slice metrics; cell-type proportion Spearman is recorded in
  the source JSON but not plotted here. Adapter color is stable across
  all three figures (\texttt{aether-3d/manuscript/compose\_palette.py}).
  \textbf{a}, Chamfer distance ($\downarrow$).
  \textbf{b}, coordinate RMSE ($\downarrow$).
  \textbf{c}, sliced 2D Wasserstein ($\downarrow$).
  \textbf{d}, Moran's I agreement on top-$K$ HVGs ($\uparrow$).
  \textbf{e}, domain ARI ($\uparrow$).
  \textbf{f}, domain NMI ($\uparrow$).
  \textbf{g}, Betti-0 stability via kNN union-find ($\uparrow$).
  \textbf{h}, direction-corrected radar over the seven plotted metrics;
  lower-is-better
  metrics are mapped via $1/(1{+}x)$ and every axis is min-max normalized
  across adapters so the outer envelope = best. One polygon per adapter
  summarizes overall reconstruction quality.
  Source: \texttt{aether-3d/results/benchmark/synthetic\_holdout.json}.}
  \label{fig:holdout}
\end{figure}

\subsection{UOT-cost component ablation}
\label{sec:res-ablation}

The $(\alpha_{\text{spatial}} \times \lambda_{\text{class}})$ grid in
Figure~\ref{fig:ablation} confirms that on synthetic data the class
penalty $\lambda_{\text{class}}$ dominates: lifting $\lambda$ from 0 to 10
raises top-1 coupling accuracy from $\sim$0.45 to $\sim$0.87--0.90 at every
$\alpha$ tested, while further increases to $\lambda = 100$ saturate. This
interpretable mechanism finding strengthens any default-justification
prose without overstating biology.

\begin{figure}[t]
  \centering
  \includegraphics[width=\textwidth]{figures/fig_uot_ablation.pdf}
  \caption{$(\alpha_{\text{spatial}} \times \lambda_{\text{class}})$ ablation
  on synthetic data.
  \textbf{a}, top-1 coupling accuracy ($\uparrow$) as a heatmap over the
  $(\alpha_{\text{spatial}}, \lambda_{\text{class}})$ grid.
  \textbf{b}, mean true-pair mass ($\uparrow$) over the same grid.
  \textbf{c}, marginal trace of top-1 accuracy along
  $\lambda_{\text{class}}$ at the central $\alpha_{\text{spatial}}$ slice
  (log-$x$).
  \textbf{d}, marginal trace along $\alpha_{\text{spatial}}$ at the
  central $\lambda_{\text{class}}$ slice. The two marginals quantify the
  slope of the response surface without requiring the reader to
  re-read the heatmap rows.
  Source: \texttt{aether-3d/results/benchmark/uot\_ablation.json}.}
  \label{fig:ablation}
\end{figure}

\subsection{Scaling}
\label{sec:res-scaling}

The scaling harness records device, torch version, CUDA version, peak
memory, runtime, and git SHA per measurement.
Figure~\ref{fig:scaling} reports runtime as a function of cells per
slice. For NumPy-only baselines (nearest-slice, linear-interp) peak GPU
memory is zero (no CUDA allocations); a future torch-using adapter wired
through the contract will produce nonzero memory under the
\texttt{dl}-style env and fail loudly on incompatible CUDA stacks.

\begin{figure}[t]
  \centering
  \includegraphics[width=\textwidth]{figures/fig_scaling_curve.pdf}
  \caption{Scaling on synthetic stacks.
  \textbf{a}, runtime (s, $\downarrow$ at fixed quality) vs cells per
  slice on log-log axes, grouped by adapter.
  \textbf{b}, per-measurement reproducibility provenance (device, torch,
  CUDA, hostname, git SHA) so a reviewer can re-run any single point.
  Source: \texttt{aether-3d/results/benchmark/scaling\_curve.json}.}
  \label{fig:scaling}
\end{figure}

\section{Discussion}
\label{sec:discussion}

\textsc{Aether3D} in its present synthetic-evaluation form is not a claim of
biological advantage over any prior 3D reconstruction method. The
audit boundary requires that every promotion of a claim from
\texttt{validated-small} or \texttt{blocked} to \texttt{paper-ready} be
backed by real-data evidence linked from the project claim ledger.
The deliverables present today are: (i)~a uniform volume-adapter
contract that catches audit regressions mechanically; (ii)~the
eight-metric per-virtual-slice suite (geometry, molecular, and
topology); (iii)~an interpretable UOT-cost
component ablation; (iv)~a scaling harness that fails
fast when the running torch cannot launch kernels on the visible GPU.

\subsection{Limitations}

Three limitations frame the present manuscript. First, the
volume-adapter contract is exercised on a synthetic point cloud with
regular inter-slice spacing; real-data sections (e.g., a BRCA IMC serial
volume) carry irregular spacing, which
shifts the burden onto the UOT cost function and the
maximal-common-region helper introduced in this revision. The
contract handles arbitrary spacing by construction, but evidence under
real organ-boundary mutations and tissue tearing is not yet in scope.
Second, the 2.5D vs continuous-3D distinction matters: our reconstruction
operates on cells as discrete coordinates rather than a continuous
volumetric field, which keeps memory predictable but means that
sub-cellular structures (vascular lumens, axonal tracts) are not
represented. Third, the vascular-network 3D-graph topology metrics ship
as a callable module
(\texttt{aether\_3d.benchmarks.graph\_3d}) but are exercised only on
synthetic graphs; the BRCA-IMC vascular-network demo~\cite{Yang2026DeepSpatial} waits
for a real IMC volume.

\subsection{Future work}

Our next step is to evaluate on a real BRCA IMC serial-section volume, at
which point the held-out virtual-slice benchmark moves from synthetic to
real evidence and the distance-band neighborhood-heatmap composer can be
re-driven against real cellular neighborhoods. Subsequent milestones add a
multi-planar (sagittal/coronal/horizontal) virtual-plane composer on real
coordinates and a per-brain-region cell-type Spearman panel on a
mouse-brain section ladder.
The whole-atlas-scale reconstruction is cited rather than reproduced~\cite{Yang2026DeepSpatial}; we plan to scale gradually,
section by section. Each
transition is gated by the claim-ledger update described above; no
claim graduates without linked evidence.
Cross-animal validation and CellCharter-style spatial-domain
analysis are likewise treated as reference-pressure checks rather than
implemented claims until the required independent atlas and external-method
contracts are present~\cite{Yang2026DeepSpatial,Varrone2024CellCharter}.

\section*{Data and code availability}

Source code, tests, and the benchmark contract live in the Aether3D
package. Synthetic placeholder data is produced deterministically by
\texttt{scripts/data/prepare\_synthetic\_aether.py}; the data card schema
and reproducibility metadata helper are documented in the package.

\bibliographystyle{plain}
\bibliography{refs}

\end{document}
